\documentclass{article}
\usepackage{graphicx}

\usepackage{amsmath, amssymb, amsthm}
\usepackage{xcolor}
\usepackage{geometry}
\geometry{left=2cm, right=2cm, top=2cm, bottom=2cm}

\newcommand{\bE}{\mathbb{E}}
\newcommand{\bQ}{\mathbb{Q}}
\newcommand{\bP}{\mathbb{P}}
\newcommand{\cF}{\mathcal{F}}

\title{MATH4907 Assignment 1}
\author{Sanyang Liu}
\date{September 2026}

\begin{document}

\maketitle

\section*{Problem 1}
Let
\[
\widetilde{V}_i:=e^{-ir\Delta t}V_i,
\qquad
T=n\Delta t,
\qquad
\Omega=\{U,D\}^n.
\]
Assume that \(S_0\) is known at time \(t_0\) and that
\(\mathcal{F}_0\) is trivial. We have
\begin{align*}
V_0
&= \widetilde{V}_0
&&\text{(by the definition of \(\widetilde{V}_i\))}\\
&= \mathbb{E}^{\mathbb{Q}}
   \left[\widetilde{V}_n\mid\mathcal{F}_0\right]
&&\text{(by the $\bQ$-martingale property)}\\
&= e^{-rT}\mathbb{E}^{\mathbb{Q}}
   \left[V_n\mid\mathcal{F}_0\right]
&&\text{(since \(\widetilde{V}_n=e^{-nr\Delta t}V_n=e^{-rT}V_n\))}\\
&= e^{-rT}\mathbb{E}^{\mathbb{Q}}
   \left[(S_n-K)^+\mid\mathcal{F}_0\right]
&&\text{(by the European call payoff)}\\
&= e^{-rT}\mathbb{E}^{\mathbb{Q}}
   \left[(S_n-K)^+\right]
&&\text{(since \(\mathcal{F}_0\) is trivial)}\\
&= e^{-rT}
   \sum_{\omega\in\Omega}
   \bigl(S_n(\omega)-K\bigr)^+
   \mathbb{Q}(\{\omega\})
&&\text{(by the definition of expectation)}\\
&= e^{-rT}
   \sum_{i=0}^{n}
   \binom{n}{i}
   q^{\,n-i}(1-q)^i
   \left(S_0u^{n-i}d^i-K\right)^+
&&\text{(grouping paths with \(i\) down moves)}\\
&= e^{-rT}
   \sum_{i=0}^{n}
   \binom{n}{i}
   q^{\,n-i}(1-q)^i
   \left(S_0u^{n-i}d^i-K\right)
   \mathbf{1}_{\{S_0u^{n-i}d^i>K\}}
&&\text{(since \(x^+=x\mathbf{1}_{\{x>0\}}\))}\\
&=
S_0
\sum_{i=0}^{n}
\binom{n}{i}
\left(uqe^{-r\Delta t}\right)^{n-i}
\left(d(1-q)e^{-r\Delta t}\right)^i
\mathbf{1}_{\{S_0u^{n-i}d^i>K\}}
\\
&\quad
-
Ke^{-rT}
\sum_{i=0}^{n}
\binom{n}{i}
q^{\,n-i}(1-q)^i
\mathbf{1}_{\{S_0u^{n-i}d^i>K\}}
&&\text{(using \(e^{-rT}=e^{-rn\Delta t}\))}\\
&=
S_0
\sum_{i=0}^{i_0}
\binom{n}{i}
\left(uqe^{-r\Delta t}\right)^{n-i}
\left(d(1-q)e^{-r\Delta t}\right)^i
\\
&\quad
-
Ke^{-rT}
\sum_{i=0}^{i_0}
\binom{n}{i}
q^{\,n-i}(1-q)^i
&&
\begin{gathered}\big(
i_0:=
\max\left\{
i\in\{0,1,\ldots,n\}:
S_0u^{n-i}d^i>K
\right\}\big)
\end{gathered}
\end{align*}
The definition of \(q\) implies
\[
uqe^{-r\Delta t}
+
d(1-q)e^{-r\Delta t}
=
e^{-r\Delta t}
\bigl(qu+(1-q)d\bigr)
=1.
\]
Therefore, defining the CDF for the binomial distribution,
\[
F^{B(n,p)}(k)
:=
\sum_{i=0}^{k}
\binom{n}{i}
p^i(1-p)^{n-i},
\]
we obtain
\[
\boxed{
V_0
=
S_0
F^{B\left(n,d(1-q)e^{-r\Delta t}\right)}(i_0)
-
Ke^{-rT}
F^{B(n,1-q)}(i_0).
}
\]

\newpage
\section*{Problem 2}
Let \(h:=\Delta t\). By the definition of \(\upsilon^2\) in (1.2) on page 22,
\[
\upsilon^2(h)=e^{(2r+\sigma^2)h}.
\]
Define
\[
A(h)
:=
\frac{1+\upsilon^2(h)}{2e^{rh}}
=
\frac12\left(e^{-rh}+e^{(r+\sigma^2)h}\right).
\]
Under the CRR condition \(ud=1\), the up and down factors are given by
\[
u=A(h)+\sqrt{A(h)^2-1},
\qquad
d=A(h)-\sqrt{A(h)^2-1}.
\]
For any fixed constant \(a\), the Taylor expansion of \(e^{ah}\) at \(h=0\) is
\[
e^{ah}=1+ah+O(h^2),
\qquad h\to0.
\]
Therefore,
\[
\begin{aligned}
A(h)
&=
\frac12\left(e^{-rh}+e^{(r+\sigma^2)h}\right)\\
&=
\frac12\left[
1-rh+O(h^2)
+
1+(r+\sigma^2)h+O(h^2)
\right]\\
&=
1+\frac12\sigma^2h+O(h^2).
\end{aligned}
\]
Consequently,
\[
\begin{aligned}
A(h)^2-1
&=
\left(1+\frac12\sigma^2h+O(h^2)\right)^2-1\\
&=
\sigma^2h+O(h^2).
\end{aligned}
\]
Assuming \(\sigma>0\), we obtain
\[
\begin{aligned}
\sqrt{A(h)^2-1}
&=
\sqrt{\sigma^2h+O(h^2)}\\
&=
\sqrt{h}\,\sqrt{\sigma^2+O(h)}\\
&=
\sigma\sqrt{h}\,\sqrt{1+O(h)}\\
&=
\sigma\sqrt{h}\left(1+O(h)\right)\\
&=
\sigma\sqrt{h}+O(h^{3/2}).
\end{aligned}
\]
Substituting these expansions into the formulas for \(u\) and \(d\), we obtain
\[
\boxed{
u
=
1+\sigma\sqrt{\Delta t}
+\frac12\sigma^2\Delta t
+O\left(\Delta t^{3/2}\right)
}
\]
and
\[
\boxed{
d
=
1-\sigma\sqrt{\Delta t}
+\frac12\sigma^2\Delta t
+O\left(\Delta t^{3/2}\right).
}
\]

\newpage
\section*{Problem 3}
Let
\[
\Delta X_i:=x_{i+1}-x_i.
\]
For the continuous model,
\[
\mathbb{E}[\Delta X_i]=\nu\Delta t,
\qquad
\operatorname{Var}(\Delta X_i)=\sigma^2\Delta t,
\]
where
\[
\nu=\mu-\frac12\sigma^2.
\]
Since
\[
\mathbb{E}[(\Delta X_i)^2]
=
\operatorname{Var}(\Delta X_i)
+
\bigl(\mathbb{E}[\Delta X_i]\bigr)^2,
\]
we have
\[
\mathbb{E}[(\Delta X_i)^2]
=
\sigma^2\Delta t+\nu^2\Delta t^2.
\]
In the trinomial tree, the increment $\Delta X_i$ takes the values
\[
\Delta x_u,\qquad 0,\qquad \Delta x_d
\]
with probabilities
\[
p_u,\qquad p_0,\qquad p_d,
\]
respectively. Matching the first moment and the second raw moment gives
\[
\begin{cases}
p_u\Delta x_u+p_d\Delta x_d=\nu\Delta t,\\[4pt]
p_u\Delta x_u^2+p_d\Delta x_d^2
=\sigma^2\Delta t+\nu^2\Delta t^2.
\end{cases}
\]
Assuming the CRR condition,
\[
ud=1
\quad\Longrightarrow\quad
\Delta x_u=-\Delta x_d=:\Delta x.
\]
Therefore,
\[
\begin{cases}
p_u\Delta x-p_d\Delta x=\nu\Delta t,\\[4pt]
p_u\Delta x^2+p_d\Delta x^2
=\sigma^2\Delta t+\nu^2\Delta t^2.
\end{cases}
\]
Dividing the two equations by $\Delta x$ and $\Delta x^2$, respectively, yields
\[
p_u-p_d=\frac{\nu\Delta t}{\Delta x},
\]
and
\[
p_u+p_d
=
\frac{\sigma^2\Delta t+\nu^2\Delta t^2}{\Delta x^2}.
\]
Adding and subtracting these two equations, we obtain
\[
\boxed{
p_u
=
\frac12\left(
\frac{\sigma^2\Delta t+\nu^2\Delta t^2}{\Delta x^2}
+
\frac{\nu\Delta t}{\Delta x}
\right)
}
\]
and
\[
\boxed{
p_d
=
\frac12\left(
\frac{\sigma^2\Delta t+\nu^2\Delta t^2}{\Delta x^2}
-
\frac{\nu\Delta t}{\Delta x}
\right).
}
\]
Finally, using
\[
p_u+p_0+p_d=1,
\]
we obtain
\[
\boxed{
p_0
=
1-
\frac{\sigma^2\Delta t+\nu^2\Delta t^2}{\Delta x^2}.
}
\]





\end{document}
